\section{Introduction}
%-------------------------- BACKGROUND ------------------------
\textbf{Background and motivation}  

Robots capable of indoor navigation hold immense utility across various sectors, including, manufacturing, logistics, healthcare, and warehouse management.
Applications involve tasks requiring movement from one location to another, such as collecting online orders from warehouses, transporting parts in assembly lines, distributing medication in hospitals, and cleaning floors in large facilities. 
The ability to navigate autonomously or semi-autonomously enables robots to perform tasks efficiently and effectively, saving time and labor costs.

Simple, lower-cost robots are more useful due to their affordability and accessibility. 
Moreover, they are less prone to malfunction due to complexity, reducing the likelihood of issues such as regular software updates, crashing, and system failures, which can hinder the reliability and effectiveness of more complex and expensive robotic systems. 

%-------------------------- DEAD RECKONING ------------------------
\textbf{Encoder odometry}

Odometry is a simple navigation method which estimates wheel rotation and translates this into changes in location and heading. 
Simple robots often rely on imperfect encoders, or even just the duration of motor activation to estimate wheel rotation. Additionally, changes in friction surfaces, changing weight loads, varying battery levels, temperature variations, wear and tear in the drive train components, and other environmental factors mean that even with perfect encoders, error will accumulate. 
As a result the navigation performance becomes compromised, hindering the ability to carry out tasks effectively in dynamic environments.

\textbf{Differential wheeled robots}
A common type of mobile robot design is a differential wheeled system, also known as two-wheeled robots or differential drive robots, which are a characterized by two independently driven wheels mounted on a common axis. These robots utilize differential steering to control their motion, where the speed and direction of each wheel can be adjusted independently to achieve different trajectories.
By varying the relative speeds of the two wheels, the robot can move forward, backward, rotate in place, or follow curved paths.
\begin{figure}[h!]
    \includegraphics[scale = 0.3]{{figures/diff_robot}}
    \caption{differential robot}
    \label{fig:acc_grid}
    \end{figure}


%-------------------------- VLP ------------------------
\textbf{visible light positioning}
Visible Light Positioning (VLP) offers significant benefits for indoor robot navigation. 
Unlike GPS, which is unreliable indoors, VLP leverages pre-existing lighting infrastructure, to provide accurate localization information. 
This makes it cost-effective and easily deployable in various indoor environments, including retail stores, warehouses, hospitals, and manufacturing facilities. 
VLP provides high-precision positioning accuracy, typically in the centimeter range, enabling robots to navigate with precision and efficiency. 
Additionally, VLP is not susceptible to signal degradation caused by obstacles or environmental factors, ensuring reliable performance even in complex indoor spaces. 
Its compatibility with existing lighting infrastructure also facilitates seamless integration into smart building systems, further enhancing its utility for indoor robot navigation applications. 

- encoding info in light

- Receiving signal  

- Extracting signal  

- Localisation algorithm  

\begin{figure}[h!]
    \includegraphics[scale = 0.1]{{figures/vlp_function}}
    \caption{grid of errors}
    \label{fig:acc_grid}
    \end{figure}

\textbf{VLP to correct odometry drift}
This paper presents a novel method of using a model with RNN architecture to combine both encoder and VLP data streams. This method combines the advantages, whilst reducing the downsides of both systems. 
It comprises of a positioning model that acts on data streams from both VLP and encoder positioning. 
The output of this model is an estimate of the current location and heading of the robot.
\begin{figure}[h!]
    \includegraphics[scale = 0.3]{{figures/nn_pre_sch}}
    \caption{The NN prediction model uses data from both VLP and }
    \label{fig:nn_pre_sch}
    \end{figure}


\textbf{Pretraining on real world VLP dataset}

\textbf{Validation in real-world environment}

\textbf{CNC}
\textbf{Mobile robot}










